\exercisesection{\S~3}

\begin{enumerate}
\def\labelenumi{\arabic{enumi})}
\item
  Let G be a compact Lie group of dimension \textgreater{} 0. Then every finite commutative subgroup of G is cyclic if and only if G is isomorphic to U, SU(2, C) or the normalizer of a maximal torus in SU(2, C).
\item
  Denote by K one of the fields R, C or H, and by n an integer \(\geq 1\) . Give the space \(K_{s}^{n}\) the usual hermitian form.
\end{enumerate}

\begin{enumerate}
\def\labelenumi{\alph{enumi})}
\item
  Show that \(\mathbf{U}(n,\mathrm{K})\) is a compact real Lie group.
\item
  Show that the sphere of radius 1 in \(\mathrm{K}_s^n\) is a (real) Lie homogeneous space for \(\mathbf{U}(n,\mathbf{K})\) ; the fixer of a point is isomorphic to \(\mathbf{U}(n - 1,\mathbf{K})\) .
\item
  Deduce that the groups \(\mathbf{U}(n,\mathbf{C})\) and \(\mathbf{U}(n,\mathbf{H})\) are connected, and that \(\mathbf{O}(n,\mathbf{R})\) has two connected components.
\item
  Show that the groups \(\mathbf{SO}(n,\mathbf{R})\) and \(\mathbf{SU}(n,\mathbf{C})\) are connected.
\item
  Show that the group \(\mathbf{O}(n,\mathbf{R})\) (resp. \(\mathbf{U}(n,\mathbf{C})\) ) is the semi-direct product of Z/2Z (resp. T) by \(\mathbf{SO}(n,\mathbf{R})\) (resp. \(\mathbf{SU}(n,\mathbf{C})\) ).
\end{enumerate}

\begin{enumerate}
\def\labelenumi{\arabic{enumi})}
\setcounter{enumi}{2}
\tightlist
\item
  \begin{enumerate}
  \def\labelenumii{\alph{enumii})}
  \tightlist
  \item
    Show that the Lie algebra of the real Lie group \(\mathbf{U}(n, \mathbf{H})\) is the set of matrices \(x \in \mathrm{M}_n(\mathbf{H})\) such that \(^t\bar{x} = -x\) , with the bracket \([x, y] = xy - yx\) . Denote it by \(\mathfrak{u}(n, \mathbf{H})\) .
  \end{enumerate}
\end{enumerate}

\begin{enumerate}
\def\labelenumi{\alph{enumi})}
\setcounter{enumi}{1}
\item
  Identify \(\mathbf{C}\) with the subfield \(\mathbf{R}(i)\) of \(\mathbf{H}\) , and \(\mathbf{C}^{2n}\) with \(\mathbf{H}^n\) by the isomorphism \((z_1, \ldots, z_{2n}) \mapsto (z_1 + jz_{n+1}, \ldots, z_n + jz_{2n})\) . Prove the equality \(\mathbf{U}(n, \mathbf{H}) = \mathbf{U}(2n, \mathbf{C}) \cap \mathbf{Sp}(2n, \mathbf{C})\) .
\item
  Deduce from \(b)\) that every compact simple (real) Lie algebra of type \(\mathbf{C}_n\) is isomorphic to \(\mathfrak{u}(n, \mathbf{H})\) .
\end{enumerate}

\begin{enumerate}
\def\labelenumi{\arabic{enumi})}
\setcounter{enumi}{3}
\tightlist
\item
  Let \(n\) be an integer \(\geq 1\) .
\end{enumerate}

\begin{enumerate}
\def\labelenumi{\alph{enumi})}
\item
  Show that the group \(\mathbf{SU}(n,\mathbf{C})\) is simply-connected (use Exerc. 2).
\item
  Show that the centre of \(\mathbf{SU}(n,\mathbf{C})\) consists of the matrices \(\lambda .I_n\) for \(\lambda \in \mathbf{C}\) , \(\lambda^n = 1\) .
\item
  Every almost simple compact Lie group of type \(A_{n}\) is isomorphic to the quotient of \(\mathbf{SU}(n+1,\mathbf{C})\) by the cyclic subgroup consisting of the matrices \(\zeta^k. I_{n+1} (0 \leq k < d)\) , where \(d\) divides \(n + 1\) and \(\zeta\) is a primitive \(d\) th root of unity.
\item
  Prove that the group \(\mathbf{SL}(n,\mathbf{C})\) is simply-connected (use Chap. III, §6, no. 9, Th. 6).
\end{enumerate}

\begin{enumerate}
\def\labelenumi{\arabic{enumi})}
\setcounter{enumi}{4}
\tightlist
\item
  For \(n\) an integer \(\geq 1\) , denote by \(\mathbf{Spin}(n, \mathbf{R})\) the reduced Clifford group associated to the usual quadratic form on \(\mathbf{R}^n\) (Algebra, Chap. IX, §9, no. 5).
\end{enumerate}

\begin{enumerate}
\def\labelenumi{\alph{enumi})}
\item
  Show that \(\mathbf{Spin}(n,\mathbf{R})\) is a compact Lie group and that the surjective homomorphism \(\varphi:\mathbf{Spin}(n,\mathbf{R})\to\mathbf{SO}(n,\mathbf{R})\) (loc. cit.) is analytic, with kernel \(\{+1,-1\}\) .
\item
  Show that, for \(n \geq 2\) , \(\mathbf{Spin}(n, \mathbf{R})\) is connected and simply-connected (use Exerc. 2). The group \(\pi_{1}(\mathbf{SO}(n, \mathbf{R}))\) is cyclic of order 2.
\item
  Let \(Z_{n}\) be the centre of \(\mathbf{Spin}(n,\mathbf{R})\) . Show that \(Z_{n}=\{+1,-1\}\) if n is odd, and \(Z_{n}=\{+1,-1,\varepsilon,-\varepsilon\}\) if n is even, where \(\varepsilon=e_{1}\ldots e_{n}\) is the product of the elements of the canonical basis of \(R^{n}\) ; the group \(Z_{2r}\) is isomorphic to \((\mathbf{Z}/2\mathbf{Z})^{2}\) (resp. to \(Z/4Z\) ) if r is even (resp. odd).
\item
  Prove that every almost simple connected compact Lie group of type \(B_{n}\) ( \(n \geq 2\) ) is isomorphic to \(\mathbf{Spin}(2n + 1, \mathbf{R})\) or \(\mathbf{SO}(2n + 1, \mathbf{R})\) .
\item
  If r is odd (resp. even) and \(\geq 2\) , every almost simple connected compact Lie group of type \(D_{r}\) is isomorphic to \(\mathbf{Spin}(2r,\mathbf{R})\) , \(\mathbf{SO}(2r,\mathbf{R})\) or \(\mathbf{SO}(2r,\mathbf{R})/\{\pm I_{2r}\}\) (resp. to one of the preceding groups or \(\mathbf{Spin}(2r,\mathbf{R})/\{1,\varepsilon\}\) ).
\end{enumerate}

\begin{enumerate}
\def\labelenumi{\arabic{enumi})}
\setcounter{enumi}{5}
\tightlist
\item
  \begin{enumerate}
  \def\labelenumii{\alph{enumii})}
  \tightlist
  \item
    Show that the compact Lie group \(\mathbf{U}(n,\mathbf{H})\) is connected and simply-connected (use Exerc. 2), and that its centre is \(\{\pm I_{n}\}\) .
  \end{enumerate}
\end{enumerate}

\begin{enumerate}
\def\labelenumi{\alph{enumi})}
\setcounter{enumi}{1}
\tightlist
\item
  Every almost simple connected compact Lie group of type \(C_{n}\) ( \(n \geq 3\) ) is isomorphic to \(\mathbf{U}(n, \mathbf{H})\) or \(\mathbf{U}(n, \mathbf{H}) / \{\pm I_{n}\}\) .
\end{enumerate}

\begin{enumerate}
\def\labelenumi{\arabic{enumi})}
\setcounter{enumi}{6}
\tightlist
\item
  Let A be the algebra of Cayley octonions (Algebra, Chap. III, App., no. 3), with the basis \((e_i)_{0\leq i\leq 7}\) of loc. cit. Denote by V the subspace of pure octonions generated by \(e_1,\ldots ,e_7\) and by E the subspace of V generated by \(e_1,e_2,e_3,e_5,e_6,e_7\) . Identify the subalgebra of A generated by \(e_0,e_4\) with the field \(\mathbf{C}\) of complex numbers, and denote by G the topological group of automorphisms of the unital algebra A.
\end{enumerate}

\begin{enumerate}
\def\labelenumi{\alph{enumi})}
\item
  Denote by Q the quadratic form on V induced by the Cayley norm, so that \((e_{i})_{1\leq i\leq7}\) is an orthonormal basis of V. Prove that the map \(\sigma\mapsto\sigma|V\) is an injective homomorphism from G to the group \(\mathbf{SO}(\mathbf{Q})\) , which is isomorphic to \(\mathbf{SO}(7,\mathbf{R})\) .
\item
  Show that the multiplication of A gives E the structure of a C-vector space of dimension 3, of which \(\{e_{1}, e_{2}, e_{3}\}\) is a basis. Denote by \(\Phi\) the hermitian form on E for which this basis is orthonormal. Let H be the fixer of \(e_{4}\) in G. Show that the map \(\sigma \mapsto \sigma|E\) is an isomorphism from H to the group \(\mathbf{SU}(\Phi)\) , which is isomorphic to \(\mathbf{SU}(3, \mathbf{C})\) . The map \(\sigma \mapsto \sigma(e_{4})\) induces an embedding of G/H in the sphere of V, which is isomorphic to \(S_{6}\) .
\item
  Let T be the torus in H consisting of the automorphisms \(\sigma\) such that \(\sigma(e_{0}) = e_{0}\) , \(\sigma(e_{1}) = \alpha e_{1}\) , \(\sigma(e_{2}) = \beta e_{2}\) , \(\sigma(e_{3}) = \gamma e_{3}\) , \(\sigma(e_{4}) = e_{4}\) , \(\sigma(e_{5}) = \bar{\alpha} e_{5}\) , \(\sigma(e_{6}) = \bar{\beta} e_{6}\) , \(\sigma(e_{7}) = \bar{\gamma} e_{7}\) , where \(\alpha, \beta, \gamma\) are three complex numbers of modulus 1 such that \(\alpha\beta\gamma = 1\) . Let N be the normalizer of T in G; show that N/T is of order 12 (note that each element of N must stabilize the set of the \(\pm e_{i}, i \neq 0\) ); deduce that G is of rank 2.
\item
  Show that G is connected semi-simple of type \(G_{2}\) and that G/H can be identified with \(S_{6}\) (show that \(G_{0} \neq H\) ; deduce that \(G_{0}\) is of type \(G_{2}\) , then use a dimension argument). Every compact group of type \(G_{2}\) is isomorphic to G.
\end{enumerate}

\begin{enumerate}
\def\labelenumi{\arabic{enumi})}
\setcounter{enumi}{7}
\item
  Let G be an almost simple, connected, compact Lie group of type \(A_{n}, B_{n}, C_{n}, D_{n}\) or \(G_{2}\) . Prove that \(\pi_{2}(G) = 0\) and \(\pi_{3}(G) = Z\) (use Exerc. 2 to 7 and the fact that \(\pi_{i}(\mathbf{S}_{n})\) is zero for i \textless{} n and cyclic for i = n (cf.~General Topology, Chap. XI).
\item
  Let g be a compact Lie algebra, t a Cartan subalgebra of g; let \((\mathrm{X}_{\alpha})_{\alpha\in\mathbb{R}}\) be a Chevalley system in the split reductive algebra \((\mathfrak{g}_{\mathbf{C}},\mathfrak{t}_{\mathbf{C}})\) such that \(X_{\alpha}\) and \(X_{-\alpha}\) are conjugate (with respect to g) for all \(\alpha\in R\) .
\end{enumerate}

\begin{enumerate}
\def\labelenumi{\alph{enumi})}
\item
  Let \(\mathcal{T}\) be a \(\mathbf{Z}\) -submodule of \(\mathfrak{t}_{\mathbf{C}}\) containing the \(iH_{\alpha}\) ( \(\alpha \in \mathrm{R}\) ) and such that \(\alpha(\mathcal{T}) \subset \mathbf{Z}i\) for all \(\alpha \in \mathrm{R}\) . Show that the \(\mathbf{Z}\) -submodule \(\mathcal{G}\) of \(\mathfrak{g}_{\mathbf{C}}\) generated by \(\mathcal{T}\) and the elements \(u_{\alpha}\) and \(v_{\alpha}\) ( \(\alpha \in \mathrm{R}\) ) is a \(\mathbf{Z}\) -Lie subalgebra of \(\mathfrak{g}\) .
\item
  Assume that g (resp. t) is the Lie algebra of a compact group G (resp. of a maximal torus T of G). Let \(\Gamma(\mathrm{T})\) be the kernel of the homomorphism \(\exp_{\mathrm{T}}:t\to\mathrm{T}\) . Show that the Z-module \((2\pi)^{-1}\varGamma(\mathrm{T})\) satisfies the hypotheses of a).
\item
  Let \(\langle,\rangle\) be an invariant scalar product on g; let \(\mu\) (resp. \(\tau\) ) be the Haar measure on g (resp. t) corresponding to the Lebesgue measure when g (resp. t) is identified with a space \(R^{n}\) by means of an orthonormal basis. Denote also by \(\mu\) (resp. \(\tau\) ) the measure on g/G (resp. t/T) that is the quotient of \(\mu\) (resp. \(\tau\) ) by the normalized Haar measure on G (resp. T). Prove the formula \(\mu(\mathfrak{g}/\mathcal{G})=\tau(\mathfrak{t}/\mathcal{T}).\prod_{\alpha\in\mathrm{R}_{+}}\langle iH_{\alpha},iH_{\alpha}\rangle\) .
\end{enumerate}

